\section{Reflections and the Staircase}

The mesh is not placed by hand. It is generated from one small chamber by mirrors. The construction that turns a short integer symbol into an infinite honeycomb is given elsewhere in the paper. This section gives the path of discovery that arrives at the symbol, the reflection-group staircase, and the exceptional cell symmetry that the climb reaches.

\begin{principle}[The crystal as a reflection orbit]
The mesh is the reflection orbit of a single chamber. A finite set of mirror walls, applied to one another forever, unfolds into all of hyperbolic four-space. A finite rule, unfolding without end.
\end{principle}

Two ideas recur and deserve a plain statement at the outset. A \emph{reflection} is the flip of a space across a flat mirror wall. It sends each point to its mirror image and fixes the wall itself. A \emph{Coxeter group} is the group generated by a small set of such reflections, with the only relations between them coming from the angles at which the walls meet. The whole of the mesh's symmetry, and so the whole of its shape, is one Coxeter group built from five mirrors.

\subsection{The atom, a single reflection}

The Coxeter group $[3,4,3,4]$ has five generating reflections. These are the five mirror walls of one fundamental chamber, the smallest tile from which the entire mesh is unfolded.

The labels $3, 4, 3, 4$ are the angles between neighbouring mirrors. Each angle is $\pi$ divided by its label.

\begin{table}[ht]
\centering
\begin{tabular}{@{}ll@{}}
\toprule
label & angle between the two mirrors \\
\midrule
$3$ & $\pi / 3$ \\
$4$ & $\pi / 4$ \\
\bottomrule
\end{tabular}
\caption{A Coxeter label is read as a mirror angle. The label $m$ between two walls means they meet at $\pi/m$.}
\label{tab:mirror-angles}
\end{table}

One chamber becomes two under a single reflection. Two become many under the rest. The mirror process recursively fills all of $H^4$.

\begin{principle}[Reflection builds, not decorates]
Repeated reflection does not decorate a space that is already there. It builds the space. The crystal is nothing more than the orbit of one chamber under its own walls.
\end{principle}

A single reflection reverses orientation, producing a mirror image. Even products of reflections preserve orientation. Odd products reverse it. This parity is the first thread of structure the mirrors hand us, and every motion of the mesh is sorted by it.

\subsection{Products of two reflections}

Two reflections composed together are the basic engine of geometry. The composition splits three ways according to how the two mirrors sit relative to each other. Crossing mirrors give a rotation. Ultraparallel mirrors, those that share no point even at infinity, give a translation. Asymptotic mirrors, those that meet only at a single boundary point at infinity, give a parabolic motion. This three-way split is the core vocabulary of the whole construction. In $\{3,4,3,4\}$ each case builds a distinct large-scale structure.

\begin{itemize}
\item \textbf{Crossing mirrors} build the local fourfold rotation. The final label $4$ in the symbol means four cells meet around each shared triangular face, and that fourfold turn is exactly the product of the two crossing walls.
\item \textbf{Ultraparallel mirrors} build a geodesic tube of cells, a one-dimensional transport spine running straight through the four-dimensional bulk.
\item \textbf{Asymptotic mirrors} build the all-important flat layer.
\end{itemize}

\subsection{The decisive case, parabolics build the flat cusp}

This is the most important construction for the substrate, so it earns its own treatment.

A \emph{parabolic} motion fixes exactly one point on the boundary at infinity. It has no fixed point anywhere inside hyperbolic space. It preserves the horospheres centred on that one boundary point, and it acts on each of them exactly like an ordinary Euclidean translation. A horosphere is the limit of a sphere whose centre has been pushed off to infinity. From inside it is intrinsically flat.

The dimension of the flat layer climbs with the dimension of the hyperbolic space that hosts it.

\begin{table}[ht]
\centering
\begin{tabular}{@{}ll@{}}
\toprule
space & parabolics build \\
\midrule
$H^2$ & flat $1$D horocycles \\
$H^3$ & flat $2$D horospheres \\
$H^4$ & flat $3$D horospheres \\
\bottomrule
\end{tabular}
\caption{The horosphere ladder. In each dimension the parabolic motions sweep out a flat layer one dimension below the bulk.}
\label{tab:horosphere-ladder}
\end{table}

Because $\{3,4,3,4\}$ lives in $H^4$, its horospheres are intrinsically flat three-dimensional spaces. This much follows from dimension alone. What makes the substrate special is the next fact.

Because the ideal vertex figure of $\{3,4,3,4\}$ is $\{4,3,4\}$, the horosphere near an ideal vertex is not merely an abstract flat three-space. It carries the cubic honeycomb. The flat layer arrives already tiled into cubes, with the spacing and symmetry fixed by the group.

\begin{result}[The flat cubic layer is the vertex figure]
In $\{3,4,3,4\}$ the flat three-dimensional cubic layer is not an extracted slice or an approximate cross section. It is the exact vertex figure at an ideal vertex, and the parabolic and affine subgroup there builds the cubic grid directly.
\end{result}

This is cleaner than the compact alternative $\{5,3,4\}$, which has no ideal vertices and no exact parabolic cusp subgroup at all. Its horospheres cut the dodecahedral cells aperiodically into a rough slab, with no clean periodic lattice on them. The flat space one finds inside $\{5,3,4\}$ is a quotient and a curvature, not an exact tiling. The flat space inside $\{3,4,3,4\}$ is a tiling, written into the symbol.

\subsection{Richer compositions, the subgroup dictionary}

Longer products of reflections give mixed motions, and each one builds a recognisable large structure.

\begin{itemize}
\item A \textbf{screw}, a rotation composed with a translation along its axis, builds a helix or chiral chain. It is richer in four dimensions, where a rotation can act in two perpendicular planes at once.
\item A \textbf{glide}, a reflection composed with a translation along the mirror, builds a staggered mirrored chain, a marching mirror.
\item A \textbf{rotary reflection} builds alternating turned-and-flipped layers.
\end{itemize}

Choosing a subgroup of $[3,4,3,4]$ chooses which large structure you build. The full group builds everything. A carefully chosen part builds one feature in isolation.

\begin{table}[ht]
\centering
\begin{tabular}{@{}lll@{}}
\toprule
subgroup & structure & dimension and type \\
\midrule
one reflection & a mirror wall & $3$D hyperbolic subspace \\
two crossing & a rotation or rosette & local, finite and cyclic \\
translation subgroup & a tube of cells & $1$D geodesic \\
screw subgroup & a helical, chiral tube & $1$D twisting \\
parabolic and affine $\{4,3,4\}$ & the cubic honeycomb layer & $3$D flat, the cusp \\
full group $[3,4,3,4]$ & the entire honeycomb & $4$D hyperbolic \\
\bottomrule
\end{tabular}
\caption{The subgroup dictionary. Each subgroup of the reflection group builds a structure of a definite dimension and kind.}
\label{tab:subgroup-dictionary}
\end{table}

The cleanest way to classify these motions is by their fixed points on the boundary at infinity, which is the three-sphere boundary of $H^4$.

\begin{table}[ht]
\centering
\begin{tabular}{@{}ll@{}}
\toprule
boundary fixed points & what it builds \\
\midrule
two & a line or tube \\
one & a flat $3$D horosphere \\
an interior fixed set & a rotation \\
\bottomrule
\end{tabular}
\caption{Motions sorted by their fixed points at infinity. The single-fixed-point parabolic case is the one that exposes flat space.}
\label{tab:boundary-fixed}
\end{table}

The one-fixed-point parabolic case is the precise mechanism by which a four-dimensional hyperbolic crystal exposes a flat three-dimensional space. It is the hinge of the whole substrate argument.

\subsection{The relations of $[3,4,3,4]$}

For a four-dimensional honeycomb the fundamental chamber is a $4$-simplex with five walls. So the group has five generators $s_0, s_1, s_2, s_3, s_4$, each its own inverse, since reflecting twice across the same wall returns every point to where it began.

What is specific to this honeycomb is the relations among the walls. For $[3,4,3,4]$ the neighbouring mirror pairs close as $(s_0 s_1)^3$, $(s_1 s_2)^4$, $(s_2 s_3)^3$, and $(s_3 s_4)^4$. All non-adjacent mirrors are perpendicular, the label $m = 2$ between them. So the whole group is captured by the diagram $\circ\,3\,\circ\,4\,\circ\,3\,\circ\,4\,\circ$.

\begin{table}[ht]
\centering
\begin{tabular}{@{}lll@{}}
\toprule
object & group & diagram \\
\midrule
$\{7,3\}$ & $[7,3]$ & $\circ\,7\,\circ\,3\,\circ$ \\
$\{5,3,4\}$ & $[5,3,4]$ & $\circ\,5\,\circ\,3\,\circ\,4\,\circ$ \\
$\{3,4,3,4\}$ & $[3,4,3,4]$ & $\circ\,3\,\circ\,4\,\circ\,3\,\circ\,4\,\circ$ \\
\bottomrule
\end{tabular}
\caption{Three hyperbolic crystals and their Coxeter diagrams. The diagram is the full specification of the symmetry.}
\label{tab:diagrams}
\end{table}

The chain of forcing is tight. The numbers fix the mirror angles. The angles fix the chamber. The chamber's reflections generate the group. The group generates the honeycomb. The honeycomb fixes the large-scale geometry. Nothing along this chain is a free choice once the symbol is named.

\subsection{From chambers to cells, the exceptional symmetry}

A single cell of the honeycomb is assembled from many chambers. The \emph{cell stabilizer}, the subgroup made of all the reflections except the final outward-pointing generator, is the cell's own internal symmetry group.

For $[3,4,3,4]$ the cell is the $24$-cell $\{3,4,3\}$, and its stabilizer is $[3,4,3]$, the full reflection group of the $24$-cell. This is the four-dimensional group associated with $F_4$, of order $1152$. It is one of the exceptional symmetry groups, with no analogue in three dimensions.

So the construction has two nested layers.

\begin{table}[ht]
\centering
\begin{tabular}{@{}lll@{}}
\toprule
layer & group & result \\
\midrule
inner & $[3,4,3]$ & chambers assemble into one $24$-cell \\
outer & $[3,4,3,4]$ & $24$-cells assemble into the honeycomb \\
\bottomrule
\end{tabular}
\caption{The two layers of the construction. The inner group builds a single cell, the outer group glues cells into the mesh.}
\label{tab:two-layers}
\end{table}

The outward generator is the wall that attaches one $24$-cell to the next across a shared octahedral facet. Dropping it leaves the cell to itself, with its full exceptional symmetry intact.

\begin{result}[The exceptional cell symmetry]
The cell stabilizer of $\{3,4,3,4\}$ is $[3,4,3]$, the order-$1152$ reflection group of $F_4$. This exceptional symmetry, with no three-dimensional counterpart, is why $\{3,4,3,4\}$ is denser and more algebraically special than the icosahedral $[5,3]$ cell of $\{5,3,4\}$.
\end{result}

The $24$-cell is the cell whose $24$ outward directions are the sites of a dock, the $24$ directed slots a tone can occupy. The exceptional symmetry of that cell is the same symmetry that gives the sites their spinor content elsewhere in the paper.

\subsection{The affine structure at infinity, where flat space hides}

The key emergence happens at an ideal vertex, a vertex pushed out to the boundary at infinity.

The vertex figure there is $\{4,3,4\}$, the cubic honeycomb. Its Coxeter group $[4,3,4]$ is Euclidean, that is, affine, rather than hyperbolic. It is the symmetry of ordinary flat three-dimensional space tiled into cubes.

\begin{principle}[A Euclidean subgroup inside the hyperbolic group]
Inside the hyperbolic group $[3,4,3,4]$ sits the Euclidean affine subgroup $[4,3,4]$. The flat cubic grid is not added to the model afterward. It is the vertex figure, already encoded in the group.
\end{principle}

\begin{table}[ht]
\centering
\begin{tabular}{@{}lll@{}}
\toprule
level & group & character \\
\midrule
the bulk & $[3,4,3,4]$ & hyperbolic \\
the flat ideal layer & $[4,3,4]$ & Euclidean, affine \\
\bottomrule
\end{tabular}
\caption{The bulk is hyperbolic, the ideal layer is flat. Both groups live inside one another, the affine one as a subgroup of the hyperbolic one.}
\label{tab:bulk-cusp}
\end{table}

This is the cleanest algebraic reason $\{3,4,3,4\}$ is special for emergent flat space. The flat three-dimensional world is not imposed on the substrate. It is read out of the substrate's own symmetry group, sitting at the corner where the hyperbolic group touches infinity.

\subsection{The conceptual staircase}

No one found $\{3,4,3,4\}$ by seeing it whole. It was reached by slowly uncovering deeper structural patterns in symmetry, regularity, reflection, curvature, and dimension. The climb runs in a definite order, and each step generalises exactly one thing.

Start with the checkerboard, the simplest repeating pattern. Move to the regular tilings of the flat plane, where only three exist, $\{3,6\}$, $\{4,4\}$, and $\{6,3\}$. Then the Platonic solids, where local angle constraints alone fix the global shape. Then recursive regularity and the Schl\"afli symbol, where a cube is written $\{4,3\}$ and the notation nests. Then higher dimensions and the regular polytopes of four-space. Then the $24$-cell, an exceptional four-dimensional shape with no three-dimensional analogue. Then reflection groups, where geometry is generated from mirrors rather than drawn by hand. Then Coxeter diagrams, where a tiny integer diagram encodes a whole geometry. Then curvature freedom, where the angle sums force the space to be spherical or hyperbolic. Then hyperbolic space itself, where the parallel postulate becomes optional and room explodes outward. Then ideal structures, vertices at infinity, horospheres, and cusps. Then the honeycombs of hyperbolic four-space. Then the cubic-honeycomb boundary, where the vertex figure $\{4,3,4\}$ of $\{3,4,3,4\}$ turns out to be flat three-space. And finally $\{3,4,3,4\}$ itself.

Each step generalised exactly one thing. Adjacency, then symmetry, then recursion, then dimension, then curvature, then ideal infinity. The order is not arbitrary. Each rung depends on the one below it.

The decisive historical milestones are few. Schl\"afli's higher-dimensional polytopes and the $24$-cell, in the eighteen-fifties. The discovery of hyperbolic geometry by Lobachevsky, Bolyai, and Gauss, in the eighteen-twenties through forties. The reflection-group viewpoint, and Coxeter's unification in the nineteen-thirties, which showed that regular polytopes, reflection groups, hyperbolic geometry, and higher dimensions are all one recursive symmetry language.

\begin{result}[The structural payoff of the climb]
The realisation that the ideal corner of $\{3,4,3,4\}$ is the ordinary cubic lattice, Euclidean three-space appearing on the boundary of hyperbolic four-space, is the payoff the whole staircase climbs toward. Flat physical space falls out of the symbol at its cusp.
\end{result}

\subsection{Why it feels discovered, not invented}

The deepest insight in the whole climb is that geometry can be generated algebraically. Not by drawing shapes by hand, but by symmetry generators, reflection relations, and combinatorial rules. By the integers alone.

These structures are rigid, highly constrained, and mathematically forced. Once the rules of regularity and curvature are accepted, very little remains free.

\begin{principle}[Forced, not chosen]
Nobody arbitrarily invented $\{3,4,3,4\}$. Once regularity and curvature are accepted, it almost reveals itself. The requirements force the answer rather than letting one pick it.
\end{principle}

This discovered-not-invented character is the theory's central claim about its own substrate. The selection does not pick $\{3,4,3,4\}$ from a menu. The requirements force it. A crystallographic spinor on the sites, negative curvature, and a flat three-dimensional cusp together leave exactly one answer.

The same logic runs through the whole theory. Local repeating rules generate large-scale order. Geometry comes from reflection. Infinite hyperbolic order emerges naturally from a handful of local laws.

\subsection{One machine, forced not invented}

The crystals $\{7,3\}$, $\{5,3,4\}$, and $\{3,4,3,4\}$ are not separate inventions. They are outputs of one construction, the reflection of a single chamber. Change only the integer symbol and the same reflection logic generates a different regular hyperbolic world. Tiny inputs, infinite outputs.

\begin{principle}[The group made visible]
The geometry is a group first. The picture is only the group made visible. The whole structure is forced by symmetry, not placed by hand, and the base stays discrete and deterministic at every step.
\end{principle}

The mesh, then, is not a drawing that happens to repeat. It is a reflection group made flesh, unfolded from five mirrors, with a flat cubic world waiting at its every cusp. The mirrors do all the work. The symbol $\{3,4,3,4\}$ is the whole of the input, and the infinite mesh of docks and sites is the whole of the output.
